% GENERATED FILE. Edit canonical lessons, then run npm run generate:books.
% canonical-source-hash: fnv1a64:2e57699df9553de9
% canonical-lessons: IT-C486-starnutire, IT-C486-prurito, IT-C486-grattarsi, IT-C486-allergia, IT-C486-puntura, IT-R486-ripasso-479-486

\chapter{Starnuti e prurito}
\label{ch:it-starnuti-e-prurito}
\hlchaptermodality{486}

\begin{chapteropening}
\textbf{By the end of this chapter:} \emph{I can describe allergies: sneezing, itching and stings.}

Complete the last lesson of chapter 486: I can describe allergies: sneezing, itching and stings.
\end{chapteropening}

\section[starnutire]{starnutire — to sneeze}

\label{lesson:IT-C486-starnutire}



\begin{quote}
Before the new one: say the Italian for gauze, then the Italian for to bandage.
\end{quote}

\subsection*{You'll want to know: starnutire}
\textbf{starnutire} — \textquotedblleft{}to sneeze\textquotedblright{}.

\begin{quote}
In primavera starnutisco tutto il giorno.
\end{quote}

\begin{quote}
\emph{In spring I sneeze all day.}
\end{quote}

An \emph{-isc-} verb. \emph{Lo starnuto} — a sneeze; people say \emph{salute!}

\begin{cousinweb}
\emph{starnutire} is inherited from Latin \emph{sternūtāre} \textquotedblleft{}to sneeze often\textquotedblright{} (Lewis \& Short), from \emph{sternuere} \textquotedblleft{}to sneeze\textquotedblright{}.
\end{cousinweb}

\begin{grammarlens}[title={What you've built}]
One more word for telling, explaining and describing at B1.
\end{grammarlens}

\subsection*{Your turn}
\begin{itemize}
\raggedright
  \item \emph{Say it:} \emph{starnutire}
  \item \emph{Say it:} \emph{starnutire}, once more
  \item \emph{Say it:} a sentence of your own with \emph{starnutire}
  \item \emph{Recall:} say the Italian for gauze, then the Italian for to bandage, then say \emph{starnutire} again
\end{itemize}

\begin{tcolorbox}[breakable,colback=teal!4,colframe=teal!35!black,title={Before you move on}]
What does starnutire mean? (\textquotedblleft{}To sneeze\textquotedblright{}.) Say it once more.
\end{tcolorbox}
\section[il prurito]{il prurito — itch, itching}

\label{lesson:IT-C486-prurito}



\begin{quote}
Before the new one: say the Italian for to bandage, then the Italian for to sneeze.
\end{quote}

\subsection*{You'll want to know: il prurito}
\textbf{il prurito} — \textquotedblleft{}itch, itching\textquotedblright{}.

\begin{quote}
Questa lana mi dà prurito.
\end{quote}

\begin{quote}
\emph{This wool makes me itch.}
\end{quote}

\emph{Avere prurito} — to itch.

\begin{cousinweb}
\emph{prurito} is a learned borrowing from Latin \emph{prūrītus} \textquotedblleft{}an itching\textquotedblright{} (Lewis \& Short), from \emph{prūrīre} \textquotedblleft{}to itch\textquotedblright{}.
\end{cousinweb}

\begin{grammarlens}[title={What you've built}]
One more word for telling, explaining and describing at B1.
\end{grammarlens}

\subsection*{Your turn}
\begin{itemize}
\raggedright
  \item \emph{Say it:} \emph{il prurito}
  \item \emph{Say it:} \emph{il prurito}, once more
  \item \emph{Say it:} a sentence of your own with \emph{il prurito}
  \item \emph{Recall:} say the Italian for to bandage, then the Italian for to sneeze, then say \emph{il prurito} again
\end{itemize}

\begin{tcolorbox}[breakable,colback=teal!4,colframe=teal!35!black,title={Before you move on}]
What does il prurito mean? (\textquotedblleft{}Itch, itching\textquotedblright{}.) Say it once more.
\end{tcolorbox}
\section[grattarsi]{grattarsi — to scratch (oneself)}

\label{lesson:IT-C486-grattarsi}



\begin{quote}
Before the new one: say the Italian for to sneeze, then the Italian for itch.
\end{quote}

\subsection*{You'll want to know: grattarsi}
\textbf{grattarsi} — \textquotedblleft{}to scratch (oneself)\textquotedblright{}.

\begin{quote}
Non grattarti, peggiora!
\end{quote}

\begin{quote}
\emph{Do not scratch, it makes it worse!}
\end{quote}

Without \emph{si}: \emph{grattare} — to scratch, to grate.

\begin{cousinweb}
\emph{grattarsi} is built on \emph{grattare} \textquotedblleft{}to scratch\textquotedblright{}, usually traced to a Germanic verb \emph{krattōn} \textquotedblleft{}to scratch\textquotedblright{}.
\end{cousinweb}

\begin{grammarlens}[title={What you've built}]
One more word for telling, explaining and describing at B1.
\end{grammarlens}

\subsection*{Your turn}
\begin{itemize}
\raggedright
  \item \emph{Say it:} \emph{grattarsi}
  \item \emph{Say it:} \emph{grattarsi}, once more
  \item \emph{Say it:} a sentence of your own with \emph{grattarsi}
  \item \emph{Recall:} say the Italian for to sneeze, then the Italian for itch, then say \emph{grattarsi} again
\end{itemize}

\begin{tcolorbox}[breakable,colback=teal!4,colframe=teal!35!black,title={Before you move on}]
What does grattarsi mean? (\textquotedblleft{}To scratch (oneself)\textquotedblright{}.) Say it once more.
\end{tcolorbox}
\section[l'allergia]{l'allergia — allergy}

\label{lesson:IT-C486-allergia}



\begin{quote}
Before the new one: say the Italian for itch, then the Italian for to scratch.
\end{quote}

\subsection*{You'll want to know: l'allergia}
\textbf{l'allergia} — \textquotedblleft{}allergy\textquotedblright{}.

\begin{quote}
Ho un'allergia ai gatti.
\end{quote}

\begin{quote}
\emph{I have an allergy to cats.}
\end{quote}

\emph{Essere allergico a} — to be allergic to.

\begin{cousinweb}
\emph{allergia} came from German \emph{Allergie} (EtymDB, under English \emph{allergy}), a modern coinage from Greek \emph{állos} \textquotedblleft{}other\textquotedblright{} and \emph{érgon} \textquotedblleft{}work, action\textquotedblright{}.
\end{cousinweb}

\begin{grammarlens}[title={What you've built}]
One more word for telling, explaining and describing at B1.
\end{grammarlens}

\subsection*{Your turn}
\begin{itemize}
\raggedright
  \item \emph{Say it:} \emph{l'allergia}
  \item \emph{Say it:} \emph{l'allergia}, once more
  \item \emph{Say it:} a sentence of your own with \emph{l'allergia}
  \item \emph{Recall:} say the Italian for itch, then the Italian for to scratch, then say \emph{l'allergia} again
\end{itemize}

\begin{tcolorbox}[breakable,colback=teal!4,colframe=teal!35!black,title={Before you move on}]
What does l'allergia mean? (\textquotedblleft{}Allergy\textquotedblright{}.) Say it once more.
\end{tcolorbox}
\section[la puntura]{la puntura — sting, bite (of an insect); injection}

\label{lesson:IT-C486-puntura}



\begin{quote}
Before the new one: say the Italian for to scratch, then the Italian for allergy.
\end{quote}

\subsection*{You'll want to know: la puntura}
\textbf{la puntura} — \textquotedblleft{}sting, bite (of an insect); injection\textquotedblright{}.

\begin{quote}
Ho una puntura di zanzara sul braccio.
\end{quote}

\begin{quote}
\emph{I have a mosquito bite on my arm.}
\end{quote}

In speech it can also be an injection: \emph{fare una puntura}.

\begin{cousinweb}
\emph{puntura} is inherited from Latin \emph{punctūra} \textquotedblleft{}a pricking\textquotedblright{} (EtymDB), from \emph{pungere} \textquotedblleft{}to prick\textquotedblright{} (Lewis \& Short).
\end{cousinweb}

\begin{grammarlens}[title={What you've built}]
One more word for telling, explaining and describing at B1.
\end{grammarlens}

\subsection*{Your turn}
\begin{itemize}
\raggedright
  \item \emph{Say it:} \emph{la puntura}
  \item \emph{Say it:} \emph{la puntura}, once more
  \item \emph{Say it:} a sentence of your own with \emph{la puntura}
  \item \emph{Recall:} say the Italian for to scratch, then the Italian for allergy, then say \emph{la puntura} again
\end{itemize}

\begin{tcolorbox}[breakable,colback=teal!4,colframe=teal!35!black,title={Before you move on}]
What does la puntura mean? (\textquotedblleft{}Sting, bite (of an insect); injection\textquotedblright{}.) Say it once more.
\end{tcolorbox}
\section[(dialogue)]{Ripasso dei capitoli 479–486}

\label{lesson:IT-R486-ripasso-479-486}



\begin{quote}
Say the Italian for cheerfulness and laugh.
\end{quote}

\subsection*{Your turn}
Say these aloud:

\begin{itemize}
\raggedright
  \item cheerfulness, laugh, to giggle, enthusiasm, to be moved
  \item to take offence, grudge, forgiveness, to regret, to let off steam
  \item hip, thigh, calf, heel, big toe
  \item eyebrow, eyelashes, eyelid, temple, jaw
  \item bruise, scratch, burn, to trip, to cut oneself
  \item X-ray, plaster cast, crutch, to put in a plaster cast, to limp
  \item to sprain, cramp, swollen, gauze, to bandage
  \item to sneeze, itch, to scratch, allergy, sting
\end{itemize}

\begin{tcolorbox}[breakable,colback=teal!4,colframe=teal!35!black,title={Before you move on}]
Cheerfulness? (\textbf{l'allegria}.) Laugh? (\textbf{la risata}.) To giggle? (\textbf{ridacchiare}.) Enthusiasm? (\textbf{l'entusiasmo}.) To be moved? (\textbf{commuoversi}.) To take offence? (\textbf{offendersi}.) Grudge? (\textbf{il rancore}.) Forgiveness? (\textbf{il perdono}.) To regret? (\textbf{pentirsi}.) To let off steam? (\textbf{sfogarsi}.) Hip? (\textbf{il fianco}.) Thigh? (\textbf{la coscia}.) Calf? (\textbf{il polpaccio}.) Heel? (\textbf{il tallone}.) Big toe? (\textbf{l'alluce}.) Eyebrow? (\textbf{il sopracciglio}.) Eyelashes? (\textbf{le ciglia}.) Eyelid? (\textbf{la palpebra}.) Temple? (\textbf{la tempia}.) Jaw? (\textbf{la mascella}.) Bruise? (\textbf{il livido}.) Scratch? (\textbf{il graffio}.) Burn? (\textbf{la scottatura}.) To trip? (\textbf{inciampare}.) To cut oneself? (\textbf{tagliarsi}.) X-ray? (\textbf{la radiografia}.) Plaster cast? (\textbf{il gesso}.) Crutch? (\textbf{la stampella}.) To put in a plaster cast? (\textbf{ingessare}.) To limp? (\textbf{zoppicare}.) To sprain? (\textbf{slogarsi}.) Cramp? (\textbf{il crampo}.) Swollen? (\textbf{gonfio}.) Gauze? (\textbf{la garza}.) To bandage? (\textbf{fasciare}.) To sneeze? (\textbf{starnutire}.) Itch? (\textbf{il prurito}.) To scratch? (\textbf{grattarsi}.) Allergy? (\textbf{l'allergia}.) Sting? (\textbf{la puntura}.)
\end{tcolorbox}
